\clearpage
\VerdierPhysicalPageAnchor{0066}%
\VerdierSemanticAnchor{ch1-2}
\section{加法圏における複体}

\VerdierSemanticAnchor{ch1-2-1}
\subsection{\texorpdfstring{$n$}{n}-重複複体}

\VerdierSemanticAnchor{ch1-2-1-1}
\noindent\textbf{2.1.1.}\quad
まず記法を定める。集合

\VerdierSemanticAnchor{ch1-2-1-1-set}
\begin{equation*}
  [n]=\{p\in\mathbb Z\mid 1\leq p\leq n\}
\end{equation*}
を $[n]$ と記す。

\VerdierSemanticAnchor{ch1-2-1-1-free-group}
集合 $[n]$ が生成する自由アーベル群を $\mathbb Z[n]$ と記し、各
$i\in[n]$ に対して、$i$ に対応する $\mathbb Z[n]$ の生成元を $e_i$ と
記す。写像 $\varphi\colon[n]\to[m]$ に対し、対応する群準同型
$\mathbb Z[n]\to\mathbb Z[m]$ も $\varphi$ と記す。この節では $C$ を
加法圏とする。

\VerdierSemanticAnchor{ch1-2-1-2}
\VerdierSemanticAnchor{ch1-2-1-2-definition}
\noindent\textbf{定義 2.1.2.}\quad
$C$ の対象からなる\emph{$n$-重複複体}とは、次のデータからなる対象である：

\begin{enumerate}
\renewcommand{\labelenumi}{\arabic{enumi})}
\item \VerdierSemanticAnchor{ch1-2-1-2-objects}%
  各 $\sigma\in\mathbb Z[n]$ に対して $C$ の対象 $X^\sigma$；
\item \VerdierSemanticAnchor{ch1-2-1-2-differentials}%
  各 $\sigma\in\mathbb Z[n]$ と各 $i\in[n]$ に対して射
  \VerdierSemanticAnchor{ch1-2-1-2-differential-map}%
  \begin{equation*}
    d^{\sigma,i}\colon X^\sigma\longrightarrow X^{\sigma+e_i},
  \end{equation*}
  ただし、これらの射は次の条件に従う：
\end{enumerate}

\VerdierSemanticAnchor{ch1-2-1-2-eq1}
\begin{equation}
  \begin{aligned}
    d^{\sigma+e_i,i}d^{\sigma,i}&=0
      &&\forall\sigma\in\mathbb Z[n],\ \forall i\in[n],\\
    d^{\sigma+e_j,i}d^{\sigma,j}
      &=d^{\sigma+e_i,j}d^{\sigma,i}
      &&\forall\sigma\in\mathbb Z[n],\ \forall i,j\in[n].
  \end{aligned}
  \tag{2.1.2.1}
\end{equation}

\VerdierSemanticAnchor{ch1-2-1-3}
\noindent\textbf{2.1.3.}\quad
$C$ の対象からなる $n$-重複複体を、通常は単に $C$ における
\emph{$n$-重複複体}と呼ぶ。$(X^\sigma,d^{\sigma,i})$ を $C$ における
$n$-重複複体とする。
\VerdierSemanticAnchor{ch1-2-1-3-component}
対象 $X^\sigma$ を複体 $(X^\sigma,d^{\sigma,i})$ の
\emph{次数 $\sigma$ の成分}と呼ぶ。
\VerdierSemanticAnchor{ch1-2-1-3-differential}
射 $d^{\sigma,i}$ を複体 $(X^\sigma,d^{\sigma,i})$ の
\emph{方向 $i$ における次数 $\sigma$ の微分}と呼ぶ。
\VerdierSemanticAnchor{ch1-2-1-3-notation}
$Y$ を $C$ における $n$-重複複体とする。$Y$ の次数 $\sigma$ の成分を
$Y^\sigma$、方向 $i$ における次数 $\sigma$ の微分を $d_Y^{\sigma,i}$ と
記す。

\VerdierSemanticAnchor{ch1-2-1-4}
\VerdierSemanticAnchor{ch1-2-1-4-definition}
\noindent\textbf{定義 2.1.4.}\quad
$Y$ と $Y'$ を $C$ における二つの $n$-重複複体とする。$Y$ から $Y'$ への
\emph{複体の射}とは、微分と可換する $C$ の射の族

\VerdierSemanticAnchor{ch1-2-1-4-family}
\begin{equation*}
  \bigl(f^\sigma\colon Y^\sigma\longrightarrow Y'^{\sigma}\bigr)_{
    \sigma\in\mathbb Z[n]}
\end{equation*}
すなわち次を満たす族である：

\VerdierSemanticAnchor{ch1-2-1-4-eq1}
\begin{equation}
  f^{\sigma+e_i}d_Y^{\sigma,i}
  =d_{Y'}^{\sigma,i}f^\sigma.
  \tag{2.1.4.1}
\end{equation}
